% ===== Chapter 7 opening: we cannot help forming groups (~5 min) =====
\partopener{Գլուխ 7 · Խմբային գրկախառնություն}{Ուրիշ մարդիկ}{Ինչպես են ուրիշ մարդիկ ազդում ուղեղի վրա (ու ինչ տարօրինակ արդյունքներ է դա տալիս)}{happy}{peach}

% The chapter's first joke: not caring only counts if people know
\begin{frame}{«Ինձ չի հետաքրքրում, թե ինչ են մտածում»}
  \begin{columns}[c]
    \begin{column}{0.46\textwidth}
      \begin{bookquote}
        Շատերը պնդում են, որ իրենց չի հետաքրքրում, թե ինչ են մտածում ուրիշները։ Ասում են դա հաճախ ու բարձրաձայն։\par\smallskip
        Կարծես արհամարհանքը չի հաշվվում, քանի դեռ այդ նույն ուրիշները դրա մասին \textbf{չեն իմացել}։
      \end{bookquote}
    \end{column}
    \begin{column}{0.5\textwidth}
      \centering
      \onslide<2->{%
      \begin{meme}[6.6cm]{mocking_spongebob.jpg}
        \memecap[text width=5.8cm]{0.5}{0.87}{Ինձ ընդհանրապես չի հետաքրքրում, թե ինչ են մտածում մարդիկ}
        \memecap[text width=5.8cm]{0.5}{0.11}{իՆձ ԸնԴհԱնՐաՊեՍ չԻ հԵտԱքՐքՐուՄ}
      \end{meme}}
    \end{column}
  \end{columns}
\end{frame}

% Rebels form groups too
\begin{frame}{Ապստամբներն էլ են խումբ կազմում}
  \begin{columns}[T]
    \begin{column}{0.31\textwidth}
      \begin{card}[colback=peach!35!cardwhite]
        \centering\bigemoji[1cm]{guitar}\par\smallskip
        {\blackfont 20-րդ դարի կես}\par
        մոդեր ու սկինհեդներ
      \end{card}
    \end{column}
    \begin{column}{0.31\textwidth}
      \onslide<2->{%
      \begin{card}[colback=lavender!35!cardwhite]
        \centering\bigemoji[1cm]{skull}\par\smallskip
        {\blackfont Այսօր}\par
        գոթեր ու էմոներ
      \end{card}}
    \end{column}
    \begin{column}{0.31\textwidth}
      \onslide<3->{%
      \begin{card}[colback=mint!35!cardwhite]
        \centering\bigemoji[1cm]{boss}\par\smallskip
        {\blackfont Միշտ}\par
        բայքերներ ու մաֆիա
      \end{card}}
    \end{column}
  \end{columns}
  \vspace{4mm}
  \onslide<4->{%
  \begin{card}[colback=butter!45!cream]
    «Ոչ ոքի նման չլինել» ուզողն առաջինը գտնում է \textbf{մեկ այլ խումբ}, որին նման լինի։ Ու բոլորը հագնվում են նույն ձևով։\par\smallskip
    Մաֆիան կարող է չհարգել օրենքը, բայց ուզում է, որ \textbf{յուրայինները} իրեն հարգեն։
  \end{card}}
\end{frame}

% Solitary confinement: contact is a need, not a wish
\begin{frame}{Մենախուցը՝ խոշտանգում}
  \begin{columns}[c]
    \begin{column}{0.54\textwidth}
      \begin{itemize}
        \item Եթե նույնիսկ կոփված հանցագործները չեն դիմանում խումբ կազմելու մղմանը, ուրեմն այն ուղեղում \textbf{շատ խորն է}
        \item<2-> Երկար ժամանակ մենախցում պահելը համարվում է \textbf{հոգեբանական խոշտանգում}
        \item<3-> Մարդկային շփումը ոչ թե ցանկություն է, այլ \textbf{անհրաժեշտություն}
      \end{itemize}
    \end{column}
    \begin{column}{0.4\textwidth}
      \centering
      \bigemoji[1.2cm]{lock}\par\medskip
      \onslide<2->{%
      \begin{card}[colback=lavender!35!cardwhite]
        ՄԱԿ-ի կանոնները (2015). մենախուց է օրական \textbf{22 ժամ} առանց իմաստալից մարդկային շփման։ \textbf{15 օրից} ավելին արգելված է։
      \end{card}}
    \end{column}
  \end{columns}
  \source{ՄԱԿ-ի «Նելսոն Մանդելայի կանոններ», 2015, կանոններ 43-44}
\end{frame}

% Other people build our brains, and their brains are... brains
\begin{frame}{Ուղեղդ կառուցել են ուրիշ ուղեղներ}
  \begin{columns}[c]
    \begin{column}{0.52\textwidth}
      \begin{itemize}
        \item Բնություն, թե դաստիարակություն (nature or nurture)։ \textbf{Երկուսն էլ}
        \item Զարգացող ուղեղի գլխավոր աղբյուրը \textbf{ուրիշ մարդիկ} են. ինչ են ասում, անում, մտածում
        \item Ինքնագնահատականն ու մոտիվացիան էլ մեծ մասամբ գալիս են ուրիշներից
      \end{itemize}
    \end{column}
    \begin{column}{0.44\textwidth}
      \centering
      \onslide<2->{%
      \begin{tikzpicture}
        \node (b) at (0,0) {\bigemoji[0.9cm]{brain}};
        \node (p) at (3.4,0) {\bigemoji[0.9cm]{people}};
        \draw[-{Stealth[length=2.6mm]}, plum, line width=1.3pt] (b.north east) to[bend left=20]
          node[above, font=\small, text=plum] {կառավարում է} (p.north west);
        \draw[-{Stealth[length=2.6mm]}, plum, line width=1.3pt] (p.south west) to[bend left=20]
          node[below, font=\small, text=plum] {ձևավորում են} (b.south east);
      \end{tikzpicture}}\par\smallskip
      \onslide<3->{%
      \begin{card}[colback=peach!35!cardwhite]
        Ուղեղները կառավարում են \textbf{իրենց իսկ զարգացումը}։ Ֆիլմերում դա անում են համակարգիչները, ու սարսափելի է։ Ուղեղներն ավելի \textbf{ծիծաղելի} են։
      \end{card}}
    \end{column}
  \end{columns}
\end{frame}
